\documentclass[11pt,a4paper]{article}

\usepackage[margin=2.5cm]{geometry}
\usepackage{booktabs}
\usepackage{tabularx}
\usepackage{array}
\usepackage{hyperref}
\usepackage{enumitem}
\usepackage{longtable}

\hypersetup{
    colorlinks=true,
    linkcolor=blue,
    urlcolor=blue
}

\title{Reproducible ANN Testing Framework\\[0.3em]
\large SIFT HDF5 Baseline on AMD Machine\\
\large (FAISS + Milvus Lite + Milvus Standalone)}
\author{Konstantin Kolchin}
\date{\today}

\begin{document}
\maketitle

\begin{abstract}
This document records a reproducible baseline workflow for ANN testing on an AMD machine,
using ready-to-use vectors (SIFT HDF5) instead of an embedding stage.
It includes setup steps, scripts to run, and measured results for FAISS, Milvus Lite,
and Milvus Standalone.
The goal is to provide a repeatable harness that any engineer can run before moving to
Milvus GPU porting work.
\end{abstract}

\tableofcontents
\clearpage

\section{Scope and Goal}

The workflow validates the first part of the demo pipeline with \textbf{ready vectors}
and avoids text chunking/embedding complexity in early iterations.
The concrete objectives are:
\begin{itemize}[leftmargin=*]
  \item We verify dataset ingest and search benchmarking on AMD host.
  \item We establish reference metrics: QPS, recall@$k$, p99 latency.
  \item We keep execution fully reproducible through explicit scripts and commands.
\end{itemize}

\section{Environment}

\begin{itemize}[leftmargin=*]
  \item Hosts: AMD machines (Hong Kong: RX 6800 XT / \texttt{gfx1030};
        7900 XTX / \texttt{gfx1100}, 2026-06-30).
  \item Working directory: \texttt{\~/ann-demo}.
  \item Python virtual environment: \texttt{.venv}.
\end{itemize}

\subsection{Python packages}

\begin{verbatim}
python3 -m venv .venv
source .venv/bin/activate
pip install --upgrade pip
pip install numpy faiss-cpu h5py
pip install "pymilvus[milvus_lite]"
\end{verbatim}

\section{Dataset Acquisition}

FTP source (\texttt{ftp.irisa.fr}) may hang on restricted networks.
The reliable path in this run was direct HTTPS with proxy bypass:

\begin{verbatim}
cd ~/ann-demo/data
wget --no-proxy https://ann-benchmarks.com/sift-128-euclidean.hdf5
\end{verbatim}

Dataset content check we use:
\begin{verbatim}
['distances', 'neighbors', 'test', 'train']
distances (10000, 100) float32
neighbors (10000, 100) int32
test      (10000, 128) float32
train   (1000000, 128) float32
\end{verbatim}

\section{Scripts and Execution}

\subsection{FAISS baseline script}
\label{sec:faiss-script}

We run:
\begin{verbatim}
python run_faiss_hdf5.py
\end{verbatim}

The script loads \texttt{train/test/neighbors} from HDF5 and runs:
\begin{itemize}[leftmargin=*]
  \item exact baseline: \texttt{IndexFlatL2},
  \item ANN baseline: \texttt{IndexIVFFlat} with variable \texttt{nprobe}.
\end{itemize}

\subsection{Milvus Lite baseline script}
\label{sec:milvus-script}

We run:
\begin{verbatim}
python run_milvus_hdf5.py
\end{verbatim}

Important implementation note: we chunk insertion to avoid gRPC message limit
(\texttt{RESOURCE\_EXHAUSTED} at 256 MB per message).

Insert pattern we recommend:
\begin{verbatim}
BATCH = 50000
for s in range(0, xb.shape[0], BATCH):
    e = min(s + BATCH, xb.shape[0])
    col.insert([ids[s:e].tolist(), xb[s:e].tolist()])
col.flush()
\end{verbatim}

\subsection{Milvus Standalone baseline script}
\label{sec:milvus-standalone-script}

We run against a Docker Compose standalone server (see \texttt{docs/QUICKSTART.md}):
\begin{verbatim}
python scripts/run_milvus_hdf5.py --uri "http://127.0.0.1:19530"
\end{verbatim}

On CPUs with SSE4.2 only (no AVX/AVX2), we use Milvus v2.5.x and set
\texttt{knowhere.simdType: sse4\_2} in \texttt{user.yaml}; otherwise index build may
exit with code 132 (\texttt{SIGILL}).

\subsection{Long runs and SSH}
\label{sec:long-runs-ssh}

Milvus benchmarks over SIFT-1M are \textbf{long-running} (often tens of minutes to
over an hour on slower hosts). If we start them over SSH without a persistent session,
a dropped connection typically \textbf{kills the job} and leaves partial or no console
output---which is easy to misread as a hang or a silent failure.

\subsubsection*{Launch pattern we recommend (copy-paste)}

\begin{verbatim}
tmux new -s milvus
cd ann-harness-amd
source .venv/bin/activate
LOG=logs/run_milvus_$(date -u +%Y%m%d_%H%M%S).log
mkdir -p logs
PYTHONUNBUFFERED=1 python scripts/run_milvus_hdf5.py 2>&1 | tee "$LOG"
# Detach: Ctrl+B, then D  (SSH can disconnect safely)
# Reattach later: tmux attach -t milvus
\end{verbatim}

For standalone Milvus, we start Docker Compose first (\texttt{milvus-docker/}), then add
\texttt{--uri "http://127.0.0.1:19530"} to the Python command. We prefer
\texttt{127.0.0.1} over \texttt{localhost} on WSL2 to avoid IPv6 client issues.

We use \texttt{PYTHONUNBUFFERED=1} (or \texttt{python -u}) so progress lines appear
immediately. Without it, stdout may be buffered when not attached to a TTY and the
run can look idle for long stretches.

\subsubsection*{Expected phases}

\begin{table}[h]
\centering
\caption{Typical Milvus harness phases (order of execution)}
\begin{tabular}{@{}p{0.22\textwidth}p{0.38\textwidth}p{0.28\textwidth}@{}}
\toprule
\textbf{Phase} & \textbf{What happens} & \textbf{Rough duration} \\
\midrule
Load HDF5 & Read \texttt{train/test/neighbors} & seconds \\
Insert 1M vectors & Chunked \texttt{col.insert}; Lite writes \texttt{milvus\_sift.db} & $\approx$1--5 min (AMD); longer on weak CPUs \\
Build IVF\_FLAT index & \texttt{create\_index} + wait & seconds--minutes (standalone often slower) \\
Load collection & \texttt{col.load()} & seconds \\
Search sweeps & Five full passes over 10k queries + p99 sampling & often the longest phase \\
\bottomrule
\end{tabular}
\end{table}

The script prints \texttt{insert\_time\_s=} and \texttt{index\_build\_time\_s=} only
\textbf{after} each phase completes; there is no progress bar during insert or search.

\subsubsection*{Monitoring while a run is active}

From a second terminal (or after \texttt{tmux attach}), we monitor with:
\begin{verbatim}
tail -f logs/run_milvus_*.log
ps -o etime,pcpu,cmd -C python          # high CPU usually means still working
du -sh milvus_sift.db                  # Lite: DB size grows during insert
docker ps --filter name=milvus         # standalone stack health
curl -sf http://127.0.0.1:9091/healthz && echo OK   # Milvus health
\end{verbatim}

\subsubsection*{Common misreadings}

\begin{itemize}[leftmargin=*]
  \item \textbf{No new log lines for 30+ minutes:} we are often still inserting, indexing, or
        searching; we check CPU and log tail before killing the process.
  \item \textbf{Same recall@10 for every \texttt{nprobe} (Milvus Lite):} the run may
        have \emph{finished successfully} while \texttt{nprobe} had no effect---not
        evidence of a crash. We use standalone Milvus for tuning sweeps.
  \item \textbf{SSH session closed:} unless the job was in \texttt{tmux}/\texttt{screen},
        we assume it stopped; we re-run inside a persistent session and archive the log file.
\end{itemize}

\section{Measured Results}

\subsection{FAISS results}

\begin{table}[h]
\centering
\caption{FAISS SIFT1M baseline (\texttt{k=10}, Hong Kong AMD host, RX 6800 XT)}
\begin{tabular}{@{}lrr@{}}
\toprule
\textbf{Configuration} & \textbf{QPS} & \textbf{recall@10} \\
\midrule
FlatL2 (exact)   & 1054.5  & 1.0000 \\
IVF, nprobe=1    & 76649.6 & 0.3707 \\
IVF, nprobe=4    & 20394.5 & 0.6984 \\
IVF, nprobe=8    & 10419.9 & 0.8336 \\
IVF, nprobe=16   & 4995.9  & 0.9274 \\
IVF, nprobe=32   & 2785.5  & 0.9761 \\
\bottomrule
\end{tabular}
\end{table}

\subsection{Milvus Lite results}

\begin{itemize}[leftmargin=*]
  \item Insert time: 79.37 s
  \item Index build time (IVF\_FLAT): 8.80 s
\end{itemize}

\begin{table}[h]
\centering
\caption{Milvus Lite IVF\_FLAT results (\texttt{k=10})}
\begin{tabular}{@{}lrrr@{}}
\toprule
\textbf{nprobe} & \textbf{QPS} & \textbf{p99 ms} & \textbf{recall@10} \\
\midrule
1  & 1490.1 & 475.36 & 0.8943 \\
4  & 1465.7 & 465.44 & 0.8943 \\
8  & 1518.8 & 456.05 & 0.8943 \\
16 & 1474.3 & 466.86 & 0.8943 \\
32 & 1535.6 & 473.83 & 0.8943 \\
\bottomrule
\end{tabular}
\end{table}

\subsection{Milvus Lite re-run (WSL2, 2026-06-26)}

Host: Intel Celeron N5095 (SSE4.2 only), Ubuntu 22.04 / WSL2,
\texttt{pymilvus 3.0.0}, \texttt{milvus-lite 3.0}.

\begin{itemize}[leftmargin=*]
  \item Insert time: 323.03 s
  \item Index build time (IVF\_FLAT): 146.56 s
\end{itemize}

\begin{table}[h]
\centering
\caption{Milvus Lite IVF\_FLAT re-run (\texttt{k=10}, WSL2)}
\begin{tabular}{@{}lrrr@{}}
\toprule
\textbf{nprobe} & \textbf{QPS} & \textbf{p99 ms} & \textbf{recall@10} \\
\midrule
1  &  262.0 & 1884.93 & 0.8942 \\
4  &  249.8 & 1679.66 & 0.8942 \\
8  &  289.0 & 1789.14 & 0.8942 \\
16 &  287.4 & 1741.21 & 0.8942 \\
32 &  279.0 & 1837.31 & 0.8942 \\
\bottomrule
\end{tabular}
\end{table}

\subsection{Milvus Standalone results (WSL2, 2026-06-26)}

Host: same WSL2 machine as above; Milvus \texttt{v2.5.4} via Docker Compose
(etcd + MinIO + standalone), URI \texttt{http://127.0.0.1:19530},
\texttt{knowhere.simdType: sse4\_2}.

\begin{itemize}[leftmargin=*]
  \item Insert time: 99.95 s
  \item Index build time (IVF\_FLAT): 457.45 s
\end{itemize}

\begin{table}[h]
\centering
\caption{Milvus Standalone IVF\_FLAT results (\texttt{k=10}, WSL2)}
\begin{tabular}{@{}lrrr@{}}
\toprule
\textbf{nprobe} & \textbf{QPS} & \textbf{p99 ms} & \textbf{recall@10} \\
\midrule
1  & 1237.2 &  27.29 & 0.3908 \\
4  &  610.7 &  33.34 & 0.7149 \\
8  &  419.6 &  47.43 & 0.8463 \\
16 &  290.2 &  34.24 & 0.9353 \\
32 &  137.7 &  57.66 & 0.9805 \\
\bottomrule
\end{tabular}
\end{table}

\subsection{FAISS results (7900 XTX host, 2026-06-30)}

Host: \texttt{amd-rx7900xtx}, Ubuntu 22.04, \texttt{faiss-cpu} via harness \texttt{.venv},
SIFT-1M HDF5 (\texttt{data/sift-128-euclidean.hdf5}), \texttt{nlist=1024}.

\begin{table}[h]
\centering
\caption{FAISS SIFT1M baseline (\texttt{k=10}, RX 7900 XTX host)}
\begin{tabular}{@{}lrr@{}}
\toprule
\textbf{Configuration} & \textbf{QPS} & \textbf{recall@10} \\
\midrule
FlatL2 (exact)   & 2340.8  & 1.0000 \\
IVF, nprobe=1    & 89993.4 & 0.3711 \\
IVF, nprobe=4    & 24484.8 & 0.6983 \\
IVF, nprobe=8    & 12510.7 & 0.8337 \\
IVF, nprobe=16   &  6408.9 & 0.9273 \\
IVF, nprobe=32   &  3298.5 & 0.9760 \\
\bottomrule
\end{tabular}
\end{table}

Recall matches Table~1 (Hong Kong) within rounding; QPS is higher on this CPU.

\subsection{Milvus Standalone via VectorDBBench (7900 XTX, 2026-06-30)}

Host: same \texttt{amd-rx7900xtx} machine; Milvus \texttt{v2.5.4} Docker Compose,
URI \texttt{http://127.0.0.1:19530}, \texttt{knowhere.simdType: sse4\_2}.
Tool: VectorDBBench \texttt{milvusivfflat}, \texttt{PerformanceCustomDataset},
SIFT-1M L2 parquet (\texttt{\textasciitilde{}vdbbench-sift1m}), \texttt{--lists 1024}.
Load + optimize ran once (\texttt{nprobe=1}); \texttt{nprobe=4..32} reused the collection
(\texttt{--skip-load}). QPS $=$ \texttt{queries}/\texttt{cost} from serial search log
(\texttt{queries=10000}).

\begin{table}[h]
\centering
\caption{VectorDBBench \texttt{milvusivfflat} SIFT-1M (\texttt{k=10}, RX 7900 XTX)}
\begin{tabular}{@{}lrrrr@{}}
\toprule
\textbf{nprobe} & \textbf{QPS} & \textbf{p99 ms} & \textbf{recall@10} & \textbf{search s} \\
\midrule
1  & 1589.0 &  0.8 & 0.3704 &  6.29 \\
4  & 1219.4 &  1.1 & 0.7017 &  8.20 \\
8  &  937.4 &  1.4 & 0.8390 & 10.67 \\
16 &  630.1 &  2.1 & 0.9312 & 15.87 \\
32 &  375.9 &  3.4 & 0.9790 & 26.60 \\
\bottomrule
\end{tabular}
\end{table}

Correct ANN curve (recall rises with \texttt{nprobe}). Search QPS is far below FAISS
(Table~5) at every step, but much higher than Milvus harness Standalone on WSL2 (Table~4).

\section{Observations}

\subsection{Per-table summary}

The measured-result tables answer different questions; we should not read them as a
single apples-to-apples shootout:
\begin{itemize}[leftmargin=*]
  \item \textbf{Table 1 (FAISS, Hong Kong / RX 6800 XT):} Reference ANN behavior. Recall rises with
        \texttt{nprobe} while QPS falls. At \texttt{nprobe=32}, recall@10 $\approx 0.98$
        at $\approx$2{,}786 QPS.
  \item \textbf{Table 2 (Milvus Lite, AMD host):} We treat this as an operational smoke test
        only. Recall is stuck at $\approx 0.89$ regardless of \texttt{nprobe}; QPS $\approx$1{,}500 with
        high p99 latency ($\approx$470 ms).
  \item \textbf{Table 3 (Milvus Lite, WSL2 re-run):} Confirms the Lite \texttt{nprobe} issue
        on a second host. Same flat recall ($\approx 0.89$), much lower QPS
        ($\approx$260--290) and higher p99 ($\approx$1.7--1.9 s).
  \item \textbf{Table 4 (Milvus Standalone harness, WSL2):} Correct ANN curve---recall@10 from 0.39
        at \texttt{nprobe=1} to 0.98 at \texttt{nprobe=32}---but search QPS is far below
        FAISS at comparable recall levels (see below).
  \item \textbf{Table 5 (FAISS, RX 7900 XTX):} Same recall ladder as Table~1 on a faster CPU;
        IVF QPS $\approx$1.18$\times$ Hong Kong at \texttt{nprobe=1} ($\approx$90k vs $\approx$77k).
  \item \textbf{Table 6 (VectorDBBench \texttt{milvusivfflat}, RX 7900 XTX):} First same-host
        Milvus IVF sweep comparable to Table~5. Recall tracks FAISS; Milvus QPS is
        $\approx$57$\times$ lower at \texttt{nprobe=1} ($\approx$1{,}589 vs $\approx$90k) but
        $\approx$3--9$\times$ higher than Table~4 at similar recall (e.g.\ $\approx$376 vs
        $\approx$138 at \texttt{nprobe=32}).
\end{itemize}

\subsection{FAISS vs Milvus Standalone at similar recall}

Pairing approximate recall targets across Table~1 and Table~4 (different hosts; see caveats):
\begin{table}[h]
\centering
\caption{FAISS (AMD) vs Milvus Standalone (WSL2) at comparable recall@10}
\begin{tabular}{@{}lrrr@{}}
\toprule
\textbf{Target recall@10} & \textbf{nprobe} & \textbf{FAISS QPS} & \textbf{Milvus QPS} \\
\midrule
$\approx 0.37$ & 1  & 76{,}650 & 1{,}237 \\
$\approx 0.83$--$0.85$ & 8  & 10{,}420 & 420 \\
$\approx 0.93$ & 16 & 4{,}996 & 290 \\
$\approx 0.98$ & 32 & 2{,}786 & 138 \\
\bottomrule
\end{tabular}
\end{table}

On raw search throughput alone, \textbf{FAISS wins decisively} in these numbers.

\subsection{FAISS vs Milvus Standalone on same host (7900 XTX)}

Tables~5 and~6 run on \texttt{amd-rx7900xtx} with the same SIFT-1M L2 dataset, IVF\_FLAT,
\texttt{nlist=1024}, and \texttt{nprobe} sweep ($1,4,8,16,32$). Recall@10 agrees within
rounding for both engines at each \texttt{nprobe} (Milvus via VectorDBBench
\texttt{milvusivfflat} on Standalone Docker CPU):
\begin{table}[h]
\centering
\caption{FAISS vs Milvus Standalone (7900 XTX): IVF \texttt{nprobe} sweep (\texttt{k=10})}
\begin{tabular}{@{}lrrr@{}}
\toprule
\textbf{nprobe} & \textbf{recall@10} & \textbf{FAISS QPS} & \textbf{Milvus QPS} \\
\midrule
1  & 0.371 & 89{,}993 & 1{,}589 \\
4  & 0.698 & 24{,}485 & 1{,}219 \\
8  & 0.834 & 12{,}511 &   937 \\
16 & 0.927 &  6{,}409 &   630 \\
32 & 0.976 &  3{,}299 &   376 \\
\bottomrule
\end{tabular}
\end{table}

On this host FAISS search QPS is $\approx$9--57$\times$ higher than Milvus at the same
\texttt{nprobe} (roughly $\approx$57$\times$ at \texttt{nprobe=1}, $\approx$9$\times$ at
\texttt{nprobe=32}). FAISS FlatL2 (exact search): QPS $=$ 2{,}341, recall@10 $=$ 1.000
(VectorDBBench run used IVF\_FLAT only).

\paragraph{How to read this gap (no sales pitch).}
Tables~5 and~6 run on the \textbf{same} host with the \textbf{same} IVF recipe. FAISS winning
by an order of magnitude (or more) is a \textbf{real} outcome---not explained away by WSL2,
SSE4.2, or ``different hardware.''

If the question is ``what should my startup use for CPU vector search?'', these numbers
point to \textbf{FAISS (or similar in-process ANN) plus the wrappers you actually need}:
a thin API service, index files on disk or object storage, metadata in Postgres, auth,
batch reindex jobs---built to your scope, without etcd/MinIO/gRPC on every query. That
architecture will stay closer to Table~5 QPS than Table~6 unless your bottleneck is
something other than search throughput (org already standardized on Milvus, multi-tenant
isolation requirements, etc.).

Listing Milvus features (filtering, RBAC, replication) does \emph{not} answer why you
should accept 9--57$\times$ lower QPS on this workload. Feature checklists are not a
substitute for the table above.

\paragraph{Why we benchmark Milvus anyway in this repo.}
\begin{itemize}[leftmargin=*]
  \item \textbf{AMD GPU goal:} hipVS/Knowhere integration targets \textbf{Milvus}, not
        \texttt{run\_faiss\_hdf5.py}. Stock FAISS CPU scripts do not exercise that path.
  \item \textbf{Cost baseline:} before GPU work, we measure what the Milvus/Knowhere \emph{CPU}
        stack adds on top of the same IVF algorithm.
  \item \textbf{External parity:} VectorDBBench and published vector-DB comparisons use
        Milvus; we need controlled CPU numbers on our hardware.
\end{itemize}
This harness does \textbf{not} argue that Milvus Standalone is the rational choice for
single-node CPU IVF search when raw QPS matters. VectorDBBench adds optimize/compact overhead
on first load; p99 latencies in Table~6 are sub-millisecond to a few ms (serial protocol).

\subsection{Does Milvus show an advantage here?}

\textbf{No---not on search QPS, including the same-host 7900 XTX comparison.} Milvus
Standalone confirms correct \texttt{nprobe} behavior (unlike Milvus Lite below) but is
much slower than in-process FAISS for this protocol. gRPC, persistence, and client
marshalling account for \emph{some} overhead; they do not make a 9--57$\times$ gap
``expected and fine'' for a team that could embed FAISS in the serving process instead.

Milvus Lite (Tables~2--3) is additionally misleading for ANN tuning: flat recall across
\texttt{nprobe} means we cannot trust sweeps. The suspiciously high Lite recall
($\approx 0.89$ vs FAISS $\approx 0.37$ at \texttt{nprobe=1}) suggests \texttt{nprobe}
is not being applied as intended.

\subsection{When a full vector DB still wins over FAISS + DIY}

Honest cases---not generic feature marketing:
\begin{itemize}[leftmargin=*]
  \item You \textbf{already need} multi-node sharding, replication, and ops tooling at
        current scale (not ``someday'').
  \item Your org or product line is committed to the Milvus/Zilliz ecosystem (support,
        managed cloud, existing integrations).
  \item You are validating \textbf{GPU search inside Milvus} on AMD (this repo's Layer~3--4
        goal)---FAISS CPU is the recall reference, not the integration target.
  \item Benchmarking and compliance require matching industry Milvus baselines, not
        fastest possible in-process search.
\end{itemize}
Otherwise, ``wrap FAISS yourself'' is a defensible founder decision given Table~5 vs
Table~6. These CPU tables omit concurrent load, HNSW/IVF\_PQ, and cluster scale-out.

FAISS strengths visible in this harness:
\begin{itemize}[leftmargin=*]
  \item Much higher QPS for the same IVF/recall trade-off on the reference AMD run.
  \item Simpler stack (no Docker, no server process).
  \item Trustworthy \texttt{nprobe} behavior in-process.
\end{itemize}

\subsection{Methodology caveats}

\begin{itemize}[leftmargin=*]
  \item Table~1 (FAISS, Hong Kong) and Tables~3--4 (Milvus harness, WSL2) run on \textbf{different hardware};
        prefer Table~5 vs Table~6 for same-host QPS comparison (7900 XTX, 2026-06-30).
  \item Table~6 uses VectorDBBench serial search (\texttt{queries=10000}); the harness uses
        the same test count but a different timing protocol and reports a p99 sample.
  \item The harness uses \texttt{IVF\_FLAT} with \texttt{nlist=1024} for both engines; other
        index types (HNSW, IVF\_PQ) may shift the picture.
  \item Milvus client code converts NumPy arrays to Python lists before search; this adds
        overhead not present in the FAISS baseline.
\end{itemize}

\subsection{Bottom line}

\begin{itemize}[leftmargin=*]
  \item Tables~1--6 do \textbf{not} show Milvus outperforming FAISS on search QPS; same-host
        proof: Table~5 vs Table~6 on 7900 XTX ($\approx$9--57$\times$ gap).
  \item For a startup optimizing CPU search latency/cost, \textbf{FAISS + minimal custom
        wrapping} is often the rational reading of these numbers---not Milvus Standalone for
        this IVF workload.
  \item We still run Milvus here as a \textbf{Knowhere/hipVS integration baseline} and for
        VectorDBBench parity, not because CPU Milvus wins on QPS.
  \item Milvus Lite is useful as a connectivity/ingest smoke test; we do not use it as an ANN tuning
        baseline.
  \item Milvus Standalone confirms proper \texttt{nprobe} behavior but is much slower than
        in-process FAISS on 7900 XTX ($\approx$376 QPS at recall@10 $\approx 0.98$ vs
        $\approx$3{,}299 FAISS; WSL2 harness Milvus $\approx$138 in Table~4).
  \item The harness goal is a \textbf{reproducible baseline before AMD GPU work}, not to
        declare Milvus the better engine for single-node CPU search.
  \item Table~5 vs Table~6 satisfies the same-host CPU comparison; next step is GPU Milvus
        (hipVS) on the same machine (Layer~1+ in \texttt{porting\_milvus\_gpu\_to\_amd.tex}).
\end{itemize}

\section{Re-run Checklist}

\begin{enumerate}[label=\textbf{R\arabic*.}, leftmargin=*]
  \item Activate environment and verify package versions.
  \item Ensure \texttt{sift-128-euclidean.hdf5} is present in \texttt{data/}.
  \item For remote hosts, start a \texttt{tmux} session (Section~\ref{sec:long-runs-ssh})
        before any Milvus run; log output with \texttt{tee}.
  \item Run \verb|python scripts/run_faiss_hdf5.py| and save console output.
  \item Run \verb|python scripts/run_milvus_hdf5.py| and save console output.
  \item Optionally run \verb|python scripts/run_milvus_hdf5.py --uri "http://127.0.0.1:19530"|
        against standalone Milvus and save console output.
  \item If insert fails with message-size error, reduce insert batch size.
  \item Archive script versions and outputs with timestamp for traceability.
\end{enumerate}

\section{Next Steps}

\begin{itemize}[leftmargin=*]
  \item \textbf{Completed on 7900 XTX lab host} (\texttt{amd-rx7900xtx}, 2026-06-30): in-repo FAISS
        harness (Table~5), Milvus Standalone Docker CPU (\texttt{v2.5.4}), and VectorDBBench
        \texttt{milvusivfflat} IVF \texttt{nprobe} sweep (Table~6;
        Section~\ref{sec:vectordbbench-ivf-nprobe}).
  \item \textbf{Layer~0 complete} on 7900 XTX lab host: ROCm \texttt{6.4.4},
        \texttt{rocminfo} shows \texttt{gfx1100}; Milvus Standalone Docker healthy
        (see \texttt{docs/porting\_milvus\_gpu\_to\_amd.tex}).
  \item \textbf{Layers~1--4} (hipRAFT/hipVS, Knowhere HIP fork, Milvus GPU build, harness GPU
        benchmarks): follow \texttt{docs/porting\_milvus\_gpu\_to\_amd.tex}.
\end{itemize}

\section{VectorDBBench (optional external benchmark)}
\label{sec:vectordbbench}

Zilliz publishes \textbf{VectorDBBench} (VDBBench), a maintained framework for comparing
vector databases: \url{https://github.com/zilliztech/vectordbbench}.
It complements---but does not replace---this harness.

\subsection{How it compares to ann-harness-amd}

\begin{table}[h]
\centering
\caption{ann-harness-amd vs VectorDBBench}
\begin{tabular}{@{}p{0.22\textwidth}p{0.34\textwidth}p{0.34\textwidth}@{}}
\toprule
\textbf{Aspect} & \textbf{ann-harness-amd} & \textbf{VectorDBBench} \\
\midrule
Purpose & Team lab notebook; AMD/Milvus bring-up & Multi-DB product benchmark + UI \\
Backends & FAISS + Milvus (Lite/Standalone) & 20+ DBs (Milvus, Qdrant, pgvector, \ldots) \\
Dataset & SIFT-1M HDF5 only & SIFT, GIST, Cohere, OpenAI, custom cases \\
Index & IVF\_FLAT, fixed \texttt{nlist/nprobe} & Per-DB (Milvus CLI: AutoIndex, FTS) \\
Best for & \texttt{nprobe} diagnostics, runbook, ROCm checks & Vendor comparison, concurrent QPS, leaderboard \\
\bottomrule
\end{tabular}
\end{table}

\begin{itemize}[leftmargin=*]
  \item We keep \textbf{ann-harness-amd} for same-host FAISS vs Milvus IVF tuning and debugging
        (e.g.\ Milvus Lite \texttt{nprobe} issues).
  \item We add \textbf{VectorDBBench} when we need official Milvus numbers, concurrent search,
        or comparison to other vector DBs (see Zilliz leaderboard:
        \url{https://zilliz.com/benchmark}).
\end{itemize}

\subsection{Workflow we recommend}

We prefer a \textbf{CLI dry-run first}, then a \textbf{serial-only} run inside \texttt{tmux}
(Section~\ref{sec:long-runs-ssh}). We skip the default concurrent search sweep on hosts with
unstable SSH.

VectorDBBench requires \textbf{Python $\geq$ 3.11} (a separate venv from this harness if
needed). Ubuntu~22.04 ships Python~3.10 only; if \texttt{python3.11} is missing, install it
first (deadsnakes PPA), then create the venv:
\begin{verbatim}
python3 --version                    # expect 3.10.x on Ubuntu 22.04
sudo apt install -y software-properties-common
sudo add-apt-repository -y ppa:deadsnakes/ppa
sudo apt update
sudo apt install -y python3.11 python3.11-venv python3.11-dev
python3.11 --version                 # must show 3.11.x before continuing

python3.11 -m venv ~/vdbbench-venv
source ~/vdbbench-venv/bin/activate
python --version                     # inside venv: 3.11.x
pip install -U pip vectordb-bench
vectordbbench --help
vectordbbench milvusautoindex --help
\end{verbatim}

If \texttt{pip install vectordb-bench} fails with ``No matching distribution'', the usual
cause is running pip under Python~3.10 (venv not activated or never created). Re-check
\texttt{python --version} inside the activated venv.

We ensure Milvus Standalone is healthy before starting. On a new host, clone this
repository first (the \texttt{milvus-docker/} stack lives inside it):
\begin{verbatim}
git clone https://github.com/konkolchin/ann-harness-amd.git
cd ann-harness-amd/milvus-docker
docker-compose up -d
docker-compose ps
curl -sf http://127.0.0.1:9091/healthz && echo OK
\end{verbatim}

On AMD hosts with SSE4.2 only, keep \texttt{user.yaml} as shipped
(\texttt{knowhere.simdType: sse4\_2}); Milvus v2.4.x may crash with \texttt{SIGILL}
without it.

\subsection{Minimal command line (local Milvus Standalone)}

VectorDBBench exposes \texttt{milvusautoindex} and \texttt{milvusfts} for Milvus;
there is no \texttt{milvusivfflat} CLI command. AutoIndex is the closest official path
but does \textbf{not} mirror our fixed IVF\_FLAT{} + \texttt{nprobe} sweep.

\textbf{Step 1: dry-run} (we validate URI and case without loading data):
\begin{verbatim}
vectordbbench milvusautoindex \
  --uri "http://127.0.0.1:19530" \
  --case-type CapacityDim128 \
  --k 10 \
  --drop-old \
  --load \
  --search-serial \
  --skip-search-concurrent \
  --db-label "local-standalone" \
  --dry-run
\end{verbatim}

\texttt{CapacityDim128} is the nearest built-in case to 128-dimensional SIFT-scale vectors.
We use \texttt{127.0.0.1} rather than \texttt{localhost} on WSL2 (IPv6 client issues).

\textbf{Step 2: real run} (after dry-run output looks correct to us):
\begin{verbatim}
tmux new -s vdbbench
LOG=logs/vdbbench_$(date -u +%Y%m%d_%H%M%S).log
mkdir -p logs
PYTHONUNBUFFERED=1 vectordbbench milvusautoindex \
  --uri "http://127.0.0.1:19530" \
  --case-type CapacityDim128 \
  --k 10 \
  --drop-old \
  --load \
  --search-serial \
  --skip-search-concurrent \
  --db-label "local-standalone" \
  2>&1 | tee "$LOG"
# Ctrl+B, then D  (detach early if SSH is unstable)
\end{verbatim}

Results are written under VectorDBBench's local results directory (path shown in dry-run
output; typically under \texttt{vectordb\_bench/results/}).

\subsection{Optional: custom SIFT-1M case}

To approximate our harness dimensions (still AutoIndex, not IVF\_FLAT):
\begin{verbatim}
vectordbbench milvusautoindex \
  --uri "http://127.0.0.1:19530" \
  --custom-case-name "SIFT1M-L2" \
  --custom-dataset-name "SIFT" \
  --custom-dataset-dim 128 \
  --custom-dataset-size 1000000 \
  --custom-dataset-metric-type L2 \
  --k 10 \
  --drop-old --load --search-serial --skip-search-concurrent \
  --dry-run
\end{verbatim}

VectorDBBench usually downloads datasets in its own layout; we confirm
\texttt{--custom-dataset-dir} and format requirements via \texttt{--help} before relying
on a local \texttt{data/sift-128-euclidean.hdf5} file.

\subsection{What not to use for this harness comparison}

\begin{itemize}[leftmargin=*]
  \item \textbf{\texttt{milvusfts}} --- full-text (BM25), not vector ANN.
  \item \textbf{\texttt{init\_bench} Web UI alone} --- we find it fine for exploration, but harder to
        archive in the runbook; we prefer CLI + \texttt{tee} for traceability.
  \item \textbf{Default concurrent search} --- much longer; does not match our serial
        QPS/recall scripts.
\end{itemize}

\subsection{Interpreting results vs Tables 1--4}

\begin{table}[h]
\centering
\caption{Expectations when comparing VectorDBBench to ann-harness-amd}
\begin{tabular}{@{}ll@{}}
\toprule
\textbf{Topic} & \textbf{Expectation} \\
\midrule
Index type & Harness: IVF\_FLAT; VDBBench Milvus CLI: AutoIndex \\
Search mode & Harness: serial batch + p99 sample; VDBBench: framework serial protocol \\
QPS / recall & Same URI does \emph{not} imply matching numbers \\
Primary use & Harness = debug/tune; VDBBench = external / multi-DB reporting \\
\bottomrule
\end{tabular}
\end{table}

\section{VectorDBBench IVF\_FLAT nprobe sweep (FAISS-comparable)}
\label{sec:vectordbbench-ivf-nprobe}

Section~\ref{sec:vectordbbench} covers VectorDBBench smoke tests (\texttt{milvusautoindex},
\texttt{CapacityDim128}, AutoIndex). \textbf{This section is separate:} we use
\texttt{milvusivfflat} to approximate the harness IVF\_FLAT{} + \texttt{nprobe} sweep
(Section~\ref{sec:milvus-script}, Tables~1--6) so we can compare recall and effective QPS
against FAISS on the same host.

\subsection{When to use this chapter}

\begin{itemize}[leftmargin=*]
  \item We want Milvus numbers from VectorDBBench (CLI logs, JSON under
        \texttt{vectordb\_bench/results/}) \emph{and} an \texttt{nprobe} ladder comparable
        to FAISS.
  \item We postpone or skip the harness \texttt{.venv} on that host but still have
        \texttt{vdbbench-venv} (Section~\ref{sec:vectordbbench}).
  \item We do \textbf{not} use this for capacity stress tests---see warnings below.
\end{itemize}

\subsection{Commands and cases (do not confuse)}

\begin{table}[h]
\centering
\caption{VectorDBBench Milvus CLI paths}
\begin{tabular}{@{}p{0.22\textwidth}p{0.36\textwidth}p{0.30\textwidth}@{}}
\toprule
\textbf{CLI} & \textbf{Index / case} & \textbf{Use for} \\
\midrule
\texttt{milvusautoindex} & AutoIndex; e.g.\ \texttt{CapacityDim128} & Smoke / Zilliz-style load; \textbf{not} \texttt{nprobe} sweep \\
\texttt{milvusivfflat} & IVF\_FLAT; \texttt{--lists}, \texttt{--probes} & \textbf{This chapter} --- \texttt{nprobe} sweep vs FAISS \\
\texttt{milvusfts} & BM25 sparse & Full-text only \\
\bottomrule
\end{tabular}
\end{table}

\textbf{CapacityDim128} repeatedly inserts SIFT batches until Milvus fails; it can fill disk
(we observed \texttt{Exit~134} and full \texttt{/} on Hong Kong). For FAISS comparison we use
\texttt{PerformanceCustomDataset} with fixed SIFT-1M parquet instead.

\subsection{Prerequisites}

Same as Section~\ref{sec:vectordbbench}: Python~$\geq$~3.11 venv, Milvus Standalone healthy
(\texttt{curl -sf http://127.0.0.1:9091/healthz}), \texttt{milvus-docker/user.yaml} with
\texttt{sse4\_2} on AMD hosts. We keep \textbf{$\geq$20--30\,GB free} on \texttt{/} before
load (insert + MinIO segments + dataset cache).

\subsection{Step 1: SIFT-1M parquet for \texttt{PerformanceCustomDataset}}

VectorDBBench has no built-in \emph{performance} case for SIFT-1M L2 (only SIFT 500K / 5M in
other case types). We convert the same ann-benchmarks HDF5 as the harness:

\begin{verbatim}
source ~/vdbbench-venv/bin/activate
pip install h5py pandas pyarrow

mkdir -p ~/ann-harness-amd/data ~/vdbbench-sift1m
wget --no-proxy -c https://ann-benchmarks.com/sift-128-euclidean.hdf5 \
  -O ~/ann-harness-amd/data/sift-128-euclidean.hdf5

python3 <<'PY'
import h5py, pandas as pd, numpy as np
from pathlib import Path
h5 = Path.home() / "ann-harness-amd/data/sift-128-euclidean.hdf5"
out = Path.home() / "vdbbench-sift1m"
out.mkdir(parents=True, exist_ok=True)
with h5py.File(h5, "r") as f:
    train, test, neighbors = f["train"][:], f["test"][:], f["neighbors"][:]
pd.DataFrame({"id": np.arange(len(train)), "emb": list(train)}).to_parquet(out / "train.parquet")
pd.DataFrame({"id": np.arange(len(test)), "emb": list(test)}).to_parquet(out / "test.parquet")
pd.DataFrame({"neighbors_id": [row.tolist() for row in neighbors]}).to_parquet(out / "neighbors.parquet")
print("wrote", out)
PY
\end{verbatim}

Dry-run one \texttt{nprobe} before a long sweep:

\begin{verbatim}
vectordbbench milvusivfflat \
  --uri "http://127.0.0.1:19530" \
  --case-type PerformanceCustomDataset \
  --custom-case-name "SIFT1M-IVF" \
  --custom-dataset-name "SIFT1M" \
  --custom-dataset-dir "$HOME/vdbbench-sift1m" \
  --custom-dataset-size 1000000 \
  --custom-dataset-dim 128 \
  --custom-dataset-metric-type L2 \
  --custom-dataset-file-count 1 \
  --lists 1024 --probes 1 --k 10 \
  --search-serial --skip-search-concurrent \
  --db-label "hk-standalone-ivf" \
  --dry-run
\end{verbatim}

\subsection{Step 2: \texttt{nprobe} sweep (load once, search five times)}

VectorDBBench accepts \textbf{one} \texttt{--probes} value per invocation (mapped to
\texttt{nprobe}). We match the harness ladder $1,4,8,16,32$ with \texttt{--lists 1024}.
First run loads data; later runs skip load and only change search params:

\begin{verbatim}
tmux new -s vdbivf
source ~/vdbbench-venv/bin/activate
cd ~/ann-harness-amd
mkdir -p logs
URI="http://127.0.0.1:19530"
LOG=logs/vdb_ivf_nprobe_$(date -u +%Y%m%d_%H%M%S).log

for NPROBE in 1 4 8 16 32; do
  echo "======== nprobe=$NPROBE ========" | tee -a "$LOG"
  if [ "$NPROBE" = "1" ]; then EXTRA="--drop-old --load"
  else EXTRA="--skip-drop-old --skip-load"; fi
  PYTHONUNBUFFERED=1 vectordbbench milvusivfflat \
    --uri "$URI" \
    --case-type PerformanceCustomDataset \
    --custom-case-name "SIFT1M-IVF" \
    --custom-dataset-name "SIFT1M" \
    --custom-dataset-dir "$HOME/vdbbench-sift1m" \
    --custom-dataset-size 1000000 \
    --custom-dataset-dim 128 \
    --custom-dataset-metric-type L2 \
    --custom-dataset-file-count 1 \
    --lists 1024 --probes "$NPROBE" --k 10 \
    --search-serial --skip-search-concurrent \
    --db-label "hk-standalone-ivf" \
    $EXTRA 2>&1 | tee -a "$LOG"
done
# Ctrl+B, then D
\end{verbatim}

\subsection{Step 3: QPS and recall from logs}

Serial search does \textbf{not} populate the summary \texttt{qps} field (often \texttt{0.0}).
We take QPS from the log line \texttt{search entire test\_data}:

\begin{verbatim}
search entire test_data: cost=1.8934s, queries=1000, avg_recall=0.8924, ...
\end{verbatim}

\textbf{Effective QPS} $= \texttt{queries} / \texttt{cost}$. Recall is \texttt{avg\_recall}
(recall@k for the case's \texttt{--k}).

\begin{verbatim}
grep -E "nprobe=|search entire test_data" logs/vdb_ivf_nprobe_*.log
\end{verbatim}

JSON results also land under
\texttt{~/vdbbench-venv/lib/python3.11/site-packages/vectordb\_bench/results/Milvus/}.

\subsection{Interpreting vs FAISS (Table~1)}

\begin{table}[h]
\centering
\caption{VDBBench IVF sweep vs harness FAISS (same host, SIFT-1M L2, \texttt{k=10})}
\begin{tabular}{@{}ll@{}}
\toprule
\textbf{Topic} & \textbf{Expectation} \\
\midrule
Index & Both IVF\_FLAT; VDBBench uses \texttt{--lists 1024} to mirror harness \texttt{nlist} \\
\texttt{nprobe} & Same ladder; VDBBench = five CLI runs, harness = one script \\
QPS & Same formula idea; VDBBench query count may differ from harness 10k test set \\
Latency & Harness reports p99 sample; VDBBench serial log gives avg/p95/p99 \\
Load path & VDBBench adds flush/compact/optimize (often tens of seconds on standalone) \\
Milvus path & Docker + gRPC vs in-process FAISS --- Milvus QPS usually lower \\
Recall curve & Should rise with \texttt{nprobe}; flat curve suggests misconfiguration \\
\bottomrule
\end{tabular}
\end{table}

Reference FAISS IVF numbers (7900 XTX host, Table~5): at \texttt{nprobe=1} recall@10
$\approx 0.37$ and QPS $\approx 9.0\times 10^4$; at \texttt{nprobe=32} recall@10
$\approx 0.98$ and QPS $\approx 3.3\times 10^3$. Measured Milvus results on the same host
are in Table~6.

\subsection{Measured results (7900 XTX, 2026-06-30)}

See Table~6 in Section ``Measured Results''; log:
\texttt{logs/vdb\_ivf\_nprobe\_*.log}.

\subsection{Results template (copy from log)}

\begin{table}[h]
\centering
\caption{VectorDBBench \texttt{milvusivfflat} SIFT-1M (template for new hosts)}
\begin{tabular}{@{}lrrr@{}}
\toprule
\textbf{nprobe} & \textbf{QPS} & \textbf{recall@10} & \textbf{queries} \\
\midrule
1  &  &  & 10000 \\
4  &  &  & 10000 \\
8  &  &  & 10000 \\
16 &  &  & 10000 \\
32 &  &  & 10000 \\
\bottomrule
\end{tabular}
\end{table}

Filled example (7900 XTX): QPS $=$ \texttt{10000}/\texttt{cost};
\texttt{nprobe=1} $\rightarrow$ 1589 QPS, recall 0.3704; \texttt{nprobe=32} $\rightarrow$
376 QPS, recall 0.9790.

\subsection{Operational notes (Hong Kong)}

\begin{itemize}[leftmargin=*]
  \item If Milvus hangs on \texttt{waiting for Milvus...}, check \texttt{docker-compose ps};
        \texttt{milvus-etcd} \texttt{Exit~1} or \texttt{milvus-standalone} \texttt{Exit~134}
        often means corrupted \texttt{volumes/} after a full disk---wipe
        \texttt{volumes/etcd}, \texttt{volumes/minio}, \texttt{volumes/milvus} together and
        \texttt{docker-compose up -d}.
  \item After a failed \texttt{CapacityDim128} run, reset volumes before IVF sweep.
  \item Keep \texttt{tmux} (\texttt{vdbivf}) for long load + five search passes over SSH.
\end{itemize}

\appendix
\section{Teammate Run Template}
\label{app:teammate-template}

This appendix provides a copy-paste template for additional benchmark runs.
We recommend one filled block per run.

\subsection*{Run Metadata}
\begin{tabularx}{\textwidth}{@{}>{\raggedright\arraybackslash}p{0.32\textwidth}X@{}}
\toprule
\textbf{Field} & \textbf{Value} \\
\midrule
Engineer & \\
Date/Time (UTC) & \\
Host name / machine ID & \\
GPU model / architecture (\texttt{gfx*}) & \\
OS / kernel & \\
ROCm / HIP version & \\
Python version & \\
Script commit hash / version tag & \\
Dataset file + checksum (if available) & \\
\bottomrule
\end{tabularx}

\subsection*{Execution Summary}
\begin{itemize}[leftmargin=*]
  \item Environment setup command set used:
  \item Dataset acquisition method used:
  \item Any deviations from this runbook:
  \item Errors/warnings observed:
\end{itemize}

\subsection*{FAISS Results (copy values)}
\begin{table}[h]
\centering
\caption{FAISS run results}
\begin{tabular}{@{}lrr@{}}
\toprule
\textbf{Configuration} & \textbf{QPS} & \textbf{recall@10} \\
\midrule
FlatL2 (exact) &  &  \\
IVF, nprobe=1  &  &  \\
IVF, nprobe=4  &  &  \\
IVF, nprobe=8  &  &  \\
IVF, nprobe=16 &  &  \\
IVF, nprobe=32 &  &  \\
\bottomrule
\end{tabular}
\end{table}

\subsection*{Milvus Results (copy values)}
\begin{itemize}[leftmargin=*]
  \item Insert time (s):
  \item Index build time (s):
\end{itemize}

\begin{table}[h]
\centering
\caption{Milvus IVF\_FLAT run results}
\begin{tabular}{@{}lrrr@{}}
\toprule
\textbf{nprobe} & \textbf{QPS} & \textbf{p99 ms} & \textbf{recall@10} \\
\midrule
1  &  &  &  \\
4  &  &  &  \\
8  &  &  &  \\
16 &  &  &  \\
32 &  &  &  \\
\bottomrule
\end{tabular}
\end{table}

\subsection*{Conclusion}
\begin{itemize}[leftmargin=*]
  \item Overall status (PASS / PARTIAL / FAIL):
  \item Primary bottleneck:
  \item Recommended next action:
\end{itemize}

\section{Script Source: \texttt{run\_faiss\_hdf5.py}}
\label{app:script-faiss}

\noindent\textbf{Last validated:} 2026-06-24 on host \texttt{Hongkong-Ubu20}
(AMD Radeon RX 6800 XT, \texttt{gfx1030}, Python 3.10, \texttt{faiss-cpu 1.14.3}).

\begin{verbatim}
import time, h5py, numpy as np, faiss

K = 10
with h5py.File("data/sift-128-euclidean.hdf5", "r") as f:
    xb = np.array(f["train"], dtype="float32")
    xq = np.array(f["test"], dtype="float32")
    gt = np.array(f["neighbors"], dtype="int64")

d = xb.shape[1]
print("xb", xb.shape, "xq", xq.shape, "gt", gt.shape)

def recall_at_k(pred, gt, k):
    return np.mean([len(set(pred[i,:k]).intersection(set(gt[i,:k]))) / k
                    for i in range(pred.shape[0])])

# Exact baseline
flat = faiss.IndexFlatL2(d)
flat.add(xb)
t0 = time.time()
_, I0 = flat.search(xq, K)
t1 = time.time()
print(f"[Flat] qps={xq.shape[0]/(t1-t0):.1f}, recall@{K}=1.0000")

# ANN baseline
nlist = 1024
ivf = faiss.IndexIVFFlat(faiss.IndexFlatL2(d), d, nlist, faiss.METRIC_L2)
ivf.train(xb[:200000])
ivf.add(xb)

for nprobe in [1, 4, 8, 16, 32]:
    ivf.nprobe = nprobe
    t0 = time.time()
    _, I = ivf.search(xq, K)
    t1 = time.time()
    r = recall_at_k(I, gt, K)
    print(f"[IVF] nprobe={nprobe:2d} qps={xq.shape[0]/(t1-t0):8.1f} recall@{K}={r:.4f}")
\end{verbatim}

\section{Script Source: \texttt{run\_milvus\_hdf5.py}}
\label{app:script-milvus}

\noindent\textbf{Last validated:} 2026-06-24 on host \texttt{Hongkong-Ubu20}
(AMD Radeon RX 6800 XT, \texttt{gfx1030}, Python 3.10, \texttt{pymilvus 3.0.0},
\texttt{milvus-lite 3.0}).

\begin{verbatim}
import time
import numpy as np
import h5py

from pymilvus import (
    connections, utility, FieldSchema, CollectionSchema, DataType, Collection
)

# ---------- config ----------
MILVUS_URI = "./milvus_sift.db"   # Milvus Lite local file
COLL = "sift1m_demo"
K = 10
NLIST = 1024
NPROBES = [1, 4, 8, 16, 32]
# ---------------------------

def recall_at_k(pred_ids, gt_ids, k):
    hits = 0
    nq = pred_ids.shape[0]
    for i in range(nq):
        hits += len(set(pred_ids[i, :k]).intersection(set(gt_ids[i, :k])))
    return hits / (nq * k)

# load hdf5
with h5py.File("data/sift-128-euclidean.hdf5", "r") as f:
    xb = np.array(f["train"], dtype=np.float32)
    xq = np.array(f["test"], dtype=np.float32)
    gt = np.array(f["neighbors"], dtype=np.int64)

d = xb.shape[1]
print(f"xb={xb.shape}, xq={xq.shape}, gt={gt.shape}, dim={d}")

# connect (Lite)
connections.connect(alias="default", uri=MILVUS_URI)

if utility.has_collection(COLL):
    utility.drop_collection(COLL)

schema = CollectionSchema(
    fields=[
        FieldSchema(name="id", dtype=DataType.INT64, is_primary=True, auto_id=False),
        FieldSchema(name="vec", dtype=DataType.FLOAT_VECTOR, dim=d),
    ],
    description="SIFT1M benchmark",
)

col = Collection(name=COLL, schema=schema)

# insert (chunked to avoid gRPC message size limit)
ids = np.arange(xb.shape[0], dtype=np.int64)
BATCH = 50000
t0 = time.time()
for s in range(0, xb.shape[0], BATCH):
    e = min(s + BATCH, xb.shape[0])
    col.insert([ids[s:e].tolist(), xb[s:e].tolist()])
col.flush()
t1 = time.time()
print(f"insert_time_s={t1 - t0:.2f}")

# build index
index_params = {
    "index_type": "IVF_FLAT",
    "metric_type": "L2",
    "params": {"nlist": NLIST},
}
t0 = time.time()
col.create_index(field_name="vec", index_params=index_params)
t1 = time.time()
print(f"index_build_time_s={t1 - t0:.2f}")

col.load()

print("\nMilvus IVF_FLAT results:")
for nprobe in NPROBES:
    search_params = {"metric_type": "L2", "params": {"nprobe": nprobe}}

    # batch search for QPS
    t0 = time.time()
    res = col.search(
        data=xq.tolist(),
        anns_field="vec",
        param=search_params,
        limit=K,
        output_fields=[],
    )
    t1 = time.time()

    total = t1 - t0
    qps = xq.shape[0] / total

    # extract ids + recall
    pred = np.array([[hit.id for hit in row] for row in res], dtype=np.int64)
    r = recall_at_k(pred, gt, K)

    # p99 latency (sampled)
    sample_n = 500
    lat_ms = []
    for i in range(sample_n):
        q = [xq[i].tolist()]
        s0 = time.time()
        _ = col.search(
            data=q,
            anns_field="vec",
            param=search_params,
            limit=K,
            output_fields=[],
        )
        s1 = time.time()
        lat_ms.append((s1 - s0) * 1000.0)
    p99 = float(np.percentile(lat_ms, 99))

    print(f"nprobe={nprobe:2d} qps={qps:8.1f} p99_ms={p99:7.2f} recall@{K}={r:.4f}")
\end{verbatim}

\end{document}
